
\begin{table}[t]
\centering
\caption{Code-level distinction among runtime variants.}
\label{tab_runtime_variant_components}
\scriptsize
\setlength{\tabcolsep}{3.4pt}
\renewcommand{\arraystretch}{1.08}
\begin{tabular}{@{}>{\raggedright\arraybackslash}p{0.18\textwidth}>{\raggedright\arraybackslash}p{0.20\textwidth}>{\raggedright\arraybackslash}p{0.20\textwidth}>{\raggedright\arraybackslash}p{0.31\textwidth}@{}}
\toprule
Component & Guard & Shield & Gated-risk \\
\midrule
Training context & Curriculum PPO & No-curriculum PPO & Curriculum PPO \\
Execution operator & TTC hard/soft cap + speed guard & Same TTC hard/soft cap + speed guard & Same TTC hard/soft cap + speed guard \\
Risk reward & No & No & Motion-gated risk reward \\
Motion/risk gate & No reward gate & No reward gate & $v_t\geq3$ km/h or $\Delta p_t>10^{-4}$; risk if TTC-risk $\geq0.01$ or lane deviation $\geq0.25$ \\
Intended role & Runtime action regularizer & Pure execution-shield baseline & Anti-stationary risk shaping plus runtime feedback \\
\bottomrule
\end{tabular}
\begin{tablenotes}
\item Guard and Shield intentionally share the same execution-side operator; they are separated to isolate the effect of curriculum/training context from the effect of runtime action intervention. Gated-risk adds a motion/risk-gated training reward on top of the same execution operator.
\end{tablenotes}
\end{table}
